%(BEGIN_QUESTION)

En dykker har lufttanktrykk på 300 bar før han går i vannet. Dykkeren går ned til 20 m dybde. Her er trykket fra vannsøylen (hydrostatisk trykk) 2 bar. Anta at luftforbruket på vei ned er så lite at det kan neglisjeres. Beregn følgende trykk i tanken:
\begin{itemize}
\item{} absolutt trykk
\item{} relativt trykk
\item{} differansetrykk (i forhold til vannet utenfor tanken)
\end{itemize}

\underbar{file i00145}
%(END_QUESTION)





%(BEGIN_ANSWER)

Absolutt trykk = 2,014.7 PSIA. Relativt trykk = 2,000 PSIG. Differansetrykk (mellom tank og vann) = 1,978 PSID.

\vskip 10pt

Relativt trykk er enkelt: det er verdien målt med manometeret (2,000 PSIG). Vi antar igjen at dykkeren ikke har redusert tanktrykket nevneverdig under nedstigning. I praksis ville tanktrykket synke litt fordi luft brukes til pusting.

Absolutt trykk er relativt trykk pluss atmosfæretrykket. Siden målingen ved vannoverflaten naturlig er ved havnivå, og atmosfæretrykket der er ca. 14.7 PSI, blir absolutt trykk i tanken 2,000 + 14.7 = 2,014.7 PSIA.

Differansetrykk er forskjellen mellom tankens relative trykk (2,000 PSI) og vannets hydrostatiske trykk (22 PSI). Det gir 1,978 PSID. Samme differanse fås også med absolutte trykk: tanken 2,014.7 PSIA og vannet 36.7 PSIA (22 + 14.7), fortsatt differanse 1,978 PSID. Nøkkelen er å bruke samme trykkreferanse hele veien: ikke bland relativt og absolutt trykk.

%(END_ANSWER)





%(BEGIN_NOTES)

%INDEX% Physics, static fluids: absolute, gauge, and differential pressures

%(END_NOTES)
